\documentclass[book.tex]{subfiles}
\begin{document}

\chapter{Learning Density Functions}


\label{chap:differentiable_dft}
\noindent\rule{11cm}{0.4pt} \\
\textbf{Prerequisites:} \autoref{chap:graphconvs}, \autoref{chap:equivariant_networks}, \autoref{chap:molecular_dynamics},  \autoref{chap:dft} \\
\textbf{Difficulty Level:} ** \\
\noindent\rule{11cm}{0.4pt}
\vspace{5mm}

Systematic learning of the exchange correlation function holds out the promise of making density functional theory methods broadly more applicable in the sciences. The core mathematical idea is that we need to construct an end-to-end differentiable implementation of the Kohn-Sham self-consistent field equations which can be back-propagated through in order to fit the exchange correlation functional from data. Learning the exchange correlational functional directly holds out the promise of applying density functional theory to a very broad range of systems beyond current state of the art \cite{kasim2021learning}, \cite{kirkpatrick2021pushing}, \cite{dick2021using}. 

%"""
%Automatic differentiation programming offers paradigm shifts to enable researchers focusing only on forward calculation without having to derive the gradients. By taking down the barrier in deriving gradients, automatic differentiation has fueled the growth of deep learning and could have the same potential in other scientific fields, such as quantum chemistry. As computational quantum chemistry involves abundant gradient calculations, writing the code with automatic differentiation allows exploration of new ideas that were prohibited by the difficulty of calculating gradients. Here we introduce an open-source differentiable quantum chemistry, DQC, and explore some applications made easy by having the automatic differentiation capability. One of the applications is learning the exchange-correlation (xc) functional in DFT. Using only eight experimental data points on diatomic molecules, our trained xc networks enable improved prediction accuracy of atomization energies across a collection of 104 molecules containing new bonds and atoms that are not present in the training dataset. Besides learning xc functional, DQC also has other applications such as learning a new basis set for a family of molecules, checking convergence stability, and predicting molecular properties with alchemical perturbation among others. To conclude the talk, I will briefly discuss the challenges in implementing software using modern automatic differentiation libraries, specifically PyTorch.
%"""

%"""
%Using an end-to-end differentiable implementation of the Kohn-Sham self-consistent field equations, we obtain an accurate neural network-based exchange and correlation (XC) functional of the electronic density. The functional is optimized using information on both energy and density while exact constraints are enforced through an appropriate neural network architecture. We evaluate our model against different families of XC approximations and show that, for non-empirical functionals, there is a strong linear correlation between energy and density errors. Using this correlation, we define a novel XC functional quality metric that includes both energy and density errors, leading to a new, improved way to rank different approximations. Judged by this metric, our machine-learned functional significantly outperforms those within the same rung of approximations.
%"""

%"""
%Density functional theory describes matter at the quantum level, but all popular approximations suffer from systematic errors that arise due to the violation of mathematical properties of the exact functional. We overcome this fundamental limitation by training a neural network on molecular data and on fictitious systems with fractional charge and spin. The resulting functional, DM21, correctly describes typical examples of artificial charge delocalization and strong correlation and performs better than traditional functionals on thorough benchmarks for main-group atoms and molecules. DM21 accurately models complex systems such as hydrogen chains, charged DNA base pairs and diradical transition states. More crucially for the field, because our methodology relies on data and constraints, which are continually improving, it represents a viable pathway towards the exact universal functional.
%"""

\section{Mathematical Basics}


Recall that the basic form of the Schrodinger's equation is given by

\begin{align}
    H \psi = E \psi
\end{align}

The Kohn-Sham equations are given by

\begin{align}
    \left ( -\frac{1}{2} \nabla^2 + v_s[n](r)\right ) \psi_i(r) = \epsilon_i \psi_i(r)
\end{align}

Here $\psi_i$ is a occupied molecular orbital and the electron density $n(r)$ is given by
\begin{align}
    n(r) &= \sum_{i=1}^N |\psi_i(r)|^2
\end{align}

An ansatz is typically applied which sets an assumption on the form of $\psi_i$. A common choice is 
\begin{align}
    \psi_i &= \sum_{\mu} C_{i\mu} \phi_\mu
\end{align}

where $\phi_\mu$ is an atom-centered Gaussian orbital. The energy can be compute in terms of the underlying electron density $n(r)$ as

\begin{align}
    E[n(r)] = \int v(r)n(r) dr + \frac{1}{2} \int \frac{n(r)n(r')}{|r - r'|^2} dr + E_{xc}[n(r)]
\end{align}

Here $E_{xc}$ is the exchange correlation density. The reformulation $E[n(r)]$ reduces the complexity from $3N$ to $3$ dimensions. The density functional theory core algorithm can be reduced to the following steps.

\begin{enumerate}
    \item Choose an orbital basis set $\{b_i(r)\}$
    \item Compute the Hamiltonian matrix $H_{ij} = \int b_i(r) \frac{\partial E}{\partial n(r)} b_j(r) dr$
    \item Do eigendecomposition of $H$ to get orbitals as eigenvectors $\psi_i(r)$
    \item Compute $\sum_i \psi(r) b(r)$ the density profile
    \item Check for convergence.
\end{enumerate}

DFT has been broadly used for materials, drug-protein interactions and finding superconducting materials. There have been multiple efforts to replace DFT calculations directly with ML for rapid approximations but the accuracy of these techniques still struggles to match classical DFT does. However, by constructing a hybrid differentiable physics model it is possible to achieve greater accuracy. We need to make all the arrows in the DFT algorithm differentiable. Here are the self consistent field iteration equations.

\begin{align}
y &= f(y, \theta) \\
\frac{\partial L}{\partial \theta} &= \frac{\partial L}{\partial y} \left ( 1 - \frac{\partial f}{\partial y} \right )^{-1} \frac{\partial f}{\partial \theta}
\end{align}

The eigendecomposition is given by
\begin{align}
\frac{\partial L}{\partial A} &= \sum_i \sum_{j \neq i} \frac{1}{\lambda_i-\lambda_j}v_i v_j^T \frac{\partial L}{\partial v_j} v_i^T \\
\frac{\partial L}{\partial A} &= - \sum_j \sum_{i \neq j, \lambda_i \neq \lambda_j} \frac{1}{\lambda_i - \lambda_j} v_iv_j^T \frac{\partial L}{\partial v_j} v_i^T
\end{align}
The second case handles degeneracy of duplicate eigenvalues. These expressions can be used to propagate back the gradient. We choose the basis set $\{b_i\}$ in terms of spherical harmonics
\begin{align}
b_i(r) = \sum_j r^{2n}c_i^{(j)} \exp(\alpha_i^{(j)} r^2) Y_{lm}(\theta, \phi)
\end{align}
$c$ and $\alpha$ are available online from basis sets databases. Analytical expressions for 
\begin{align}
H_{ij} = \int b_i(r) \frac{\partial E}{\partial n(r)} b_j(r) dr 
\end{align}
are already available. Derivatives $\frac{\partial b}{\partial c}$ take the same form.

A number of these gradients are challenging for automatic differentiation systems to compute. The full hybrid neural network functional ius given by

\begin{align}
    E_{nnLDA}[n] = \alpha E_{LDA(xc)}[n] + \beta \int n(r) f(n, \xi) dr \\
    E_{nnPBE}[n] = \alpha E_{PBA(xc)}[n] + \beta \int n(r) f(n, \xi, s) dr
\end{align}
where $f$ is a fully connected neural network.

\section{Choosing the Loss Function}

A differentiable physics DFT models can be trained by comparing to a classical density functional theoretic code. We define the loss as the $L^2$ loss against energies computed by a reference code. This code can be another DFT code or possibly computed by a higher level of quantum calculation

\begin{align}
\mathcal{L} &= \sum_{i=1}^M |E_{\theta}(x) - E_{DFT}(x)|^2
\end{align}

%\textcolor{red}{TODO: Does this mean that a differentiable physics DFT equation can only be as accurate as the underlying base DFT code?}

%\section{Applying Constraints}

%\textcolor{red}{Notes: Deepmind applies two constraints to ensure correctness of fractional charge and fractional spin.}

\section{Discussion}
In principle, the learned differentiable exchange correlation functional can be arbitrarily accurate. Past research reports that learned functions can be more accurate than CCSD calculations in some cases even. % But less than CCSD(T) or SCAN. 

Applying an appropriate density regularization scheme appears to be a critical idea for good performance. Considerable additional work will need to be done though before differentiable exchange correlation functionals can be applied successfully for practical applications.

%Interested in learning the basis set in quantum chemistry calculations. In some early results, see some large improvements on the basis set calculations.

%Also interested in alchemical perturbation of diatomic molecules.
\end{document}